%%%%%%%%%%%%%%%%%%%%%%%%%%%%%%%%%%%%%%%%%%%%%%%%%%%%%%%
% Dados do Trabalho de Conclusão de Curso
%%%%%%%%%%%%%%%%%%%%%%%%%%%%%%%%%%%%%%%%%%%%%%%%%%%%%%%
\tipotrabalho{Trabalho de Conclusão de Curso II}

\providecommand{\terCapa}{Sim}
\providecommand{\terFolhaRosto}{Sim}
\providecommand{\terTermoAprovacao}{Nao}
\providecommand{\terDedicatoria}{Nao}
\providecommand{\terFichaCatalografica}{Nao}
\providecommand{\terEpigrafe}{Nao}
\providecommand{\terAgradecimentos}{Nao}
\providecommand{\terErrata}{Nao}
\providecommand{\terListaFiguras}{Sim}
\providecommand{\terListaQuadros}{Nao}
\providecommand{\terListaTabelas}{Sim}
\providecommand{\terSiglasAbrev}{Sim}
\providecommand{\terSimbolos}{Sim}
\providecommand{\terResumos}{Sim}
\providecommand{\terSumario}{Sim}
\providecommand{\terAnexo}{Nao}
\providecommand{\terApendice}{Sim}
\providecommand{\terIndiceR}{Nao}

% Título original em inglês: Pattern Analysis for Anomaly Localization in Distribution Networks with Machine Learning Approaches.
\titulo{Análise de padrões para localização de anomalias em redes de distribuição com abordagens de aprendizado de máquina}
\autor{Matheus Mion Chollet}
\local{Curitiba}
\data{2026}

% Titulação e vínculo institucional do orientador não informados.
\orientador{Leandro dos Santos Coelho}
\orientadora{}
\coorientador{}
\coorientadora{}
\scoorientador{}
\scoorientadora{}

\addbibresource{referencias.bib}

\instituicao{Universidade Federal do Paraná}
\def \ImprimirSetor{Setor de Tecnologia}
\def \ImprimirProgramaPos{}
\def \ImprimirCurso{Curso de Engenharia Elétrica}

\preambulo{Trabalho de Conclusão de Curso II apresentado ao Curso de Engenharia Elétrica do Setor de Tecnologia da Universidade Federal do Paraná como requisito parcial à obtenção do título de Engenheiro Eletricista}

\newcommand{\imprimirCurso}{Curso de Engenharia Elétrica}
% Ano de defesa adotado: 2026, conforme o ano acadêmico do projeto.
\newcommand{\imprimirDataDefesa}{04/12/2026}
\newcommand{\imprimircdu}{}
\newcommand{\imprimirerrata}{}

\providecommand{\imprimirdataapresentacaoRotulo}{}
\providecommand{\imprimirdataapresentacao}{}
\newcommand{\dataapresentacao}[2][\dataapresentacaoname]{\renewcommand{\dataapresentacao}{#2}}
\providecommand{\imprimirnomedocursoRotulo}{}
\providecommand{\imprimirnomedocurso}{}
\newcommand{\nomedocurso}[2][\nomedocursoname]{\renewcommand{\imprimirnomedocursoRotulo}{#1}\renewcommand{\imprimirnomedocurso}{#2}}
% Banca informada pelo autor; titulações e instituições a confirmar.
% O termo de aprovação permanece desativado até o recebimento do documento oficial.
\newcommand{\BancaMembroUm}{Marcos Vinicio Haas Rambo}
\newcommand{\BancaMembroDois}{Marlio José do Couto Bonfim}
\newcommand{\BancaMembroTres}{Leandro dos Santos Coelho}
\newcommand{\AssinaAprovacao}{%
  \assinatura{\BancaMembroUm}%
  \assinatura{\BancaMembroDois}%
}
% Restrições: preservar o sigilo de dados identificáveis das unidades consumidoras
% e de informações operacionais da empresa detentora dos dados.
% Ficha catalográfica: aplicabilidade a confirmar; não gerar dados fictícios.
\newcommand{\EpigrafeTexto}{}

\newcommand{\ResumoTexto}{%
A medição inteligente permite observar consumo, reversão de fluxo e disponibilidade da comunicação em escalas inferiores ao ciclo de faturamento. Desenvolveu-se uma plataforma para detectar, localizar e caracterizar anomalias em redes de distribuição, com ênfase em padrões compatíveis com perdas não técnicas. O pipeline reconstruiu por vigência o vínculo entre unidade consumidora e medidor, saneou fisicamente desequilíbrios sob carga e cadências heterogêneas, impôs isolamento programático do subespaço elétrico e combinou seis detectores não supervisionados com normalização quantílica congelada, detecção de mudança de regime (CUSUM), perfis intradiários de 24h, alarmes físicos, roteamento hierárquico por alimentador (MoE) e contraste espacial para telecomunicações. A avaliação computacional abrangeu análise exploratória, inferência fora do período de ajuste em dois cortes (19/11 e 25/11/2026), auditorias de sensibilidade com modelo congelado, ablação sob orçamento fixo e 23 cenários de perturbações fisicamente coerentes sobre 100 séries reais. Em 19/11/2026, avaliaram-se 65.686 unidades (62.857 elegíveis e 6.037 candidatas na referência \texttt{v1}; 4.298 na variante \texttt{v2} saneada, com 1.377 acionamentos de $S_{tamper}$). A carteira Top-2.000 mostrou robustez a variações de $\pm 20\%$ nos pesos (Jaccard $\ge 0{,}972$) e persistência de 88\% a 91,5\% em seis dias adicionais reais. Sob orçamento fixo, a arquitetura multivisão zerou o vazamento de quarentenas de comunicação no Top-2.000, elevou o F1 controlado de 68,25\% para 76,00\%, localizou quedas de $-50\%$ com erro mediano de 0,0~dia e isolou 16/16 eventos coletivos de circuito. A confirmação de irregularidades e o rendimento de inspeções requerem verificação de campo posterior.}

\newcommand{\PalavraschaveTexto}{AMI; detecção de anomalias; perdas não técnicas; avaliação retrospectiva; séries temporais; aprendizado de máquina.}

\newcommand{\AbstractTexto}{%
Smart metering makes consumption, reverse energy flow and communication availability observable at time scales shorter than the billing cycle. A platform was developed to detect, localize and characterize anomalies in distribution networks, focusing on patterns consistent with non-technical losses. The pipeline reconstructed time-valid links between service points and physical meters, physically sanitized loaded current imbalance and heterogeneous cadences, enforced programmatic isolation of the electrical subspace, and combined six unsupervised detectors with frozen quantile normalization, change-point detection (CUSUM), 24h intraday profiles, physical alarms, feeder-level Mixture of Experts routing, and spatial telecom contrast. Computational evaluation comprised exploratory analysis, out-of-time inference across two cutoffs (November 19 and 25, 2026), frozen-model sensitivity audits, fixed-budget ablation, and 23 physically coherent perturbation scenarios over 100 real AMI series. On November 19, 2026, 65,686 service points were evaluated (62,857 eligible and 6,037 candidates in baseline \texttt{v1}; 4,298 in hardened \texttt{v2}, with 1,377 $S_{tamper}$ triggers). The Top-2,000 queue remained robust under $\pm 20\%$ weight perturbations (Jaccard $\ge 0.972$) and showed 88.0\% to 91.5\% persistence across six additional real days. Under fixed budget, the multiview architecture eliminated telemetry-quarantine leakage in the Top-2,000, raised controlled electrical F1 from 68.25\% to 76.00\%, localized $-50\%$ drops with 0.0-day median error, and isolated 16/16 collective circuit events. Confirmation of irregularities and inspection yield require subsequent field verification.}

\newcommand{\KeywordsTexto}{AMI; anomaly detection; non-technical losses; retrospective evaluation; time series; machine learning.}
\newcommand{\Resume}{}
\newcommand{\Motscles}{}
\newcommand{\Resumen}{}
\newcommand{\Palabrasclave}{}
\newcommand{\AgradecimentosTexto}{}
\newcommand{\DedicatoriaTexto}{}
